% Class-compatibility preamble retained for a reader excerpt of
%   Sahar Diskin, Anna Geisler, "Perfect Matching in Product Graphs and in their
%   Random Subgraphs", arXiv:2404.14020v4, CC BY 4.0.
%
% The source pins its own publisher class (scrartcl); the delivery wrapper accepts
% only the classes article and amsart, so this file supplies the amsart baseline and
% then copies the author's own declaration block mechanically and unchanged
% (pinned-original-source.tex, source lines 30-57: the author's newtheorem
% declarations and the author's macro declarations).  The author's own package
% list of source lines 4-21 is reused only for the entries the retained text needs
% (enumitem, amsmath, amssymb, amsthm, mathtools, xcolor, geometry) and the
% cross-reference packages already needed by the author's own \Cref/\eqref usage
% (hyperref, cleveref); tikz, caption, subcaption, float, graphicx, wrapfig and
% abstract serve figures and the title block that this excerpt does not carry.
% No mathematical macro, statement, symbol or assumption is introduced, changed
% or removed by this file.
\documentclass[10pt,a4paper]{amsart}
\usepackage[shortlabels]{enumitem}
\usepackage{amsmath}
\usepackage{amssymb}
\usepackage{amsthm}
\usepackage{mathtools}
\usepackage{xcolor}
\usepackage[margin=1in]{geometry}
\usepackage{hyperref}
\usepackage{cleveref}
\hypersetup{
  colorlinks=true,
  linkcolor=blue,
  filecolor=magenta,
  urlcolor=cyan
}

\newtheorem{theorem}{Theorem}
\newtheorem{thm}{Theorem}[section]
\newtheorem{corollary}[thm]{Corollary}
\newtheorem{lemma}[thm]{Lemma}
\newtheorem{proposition}[thm]{Proposition}
\newtheorem{claim}[thm]{Claim}
\newtheorem{remark}[thm]{Remark}
\newtheorem{question}[thm]{Question}

\newcommand{\whp}{whp}
\newcommand{\defi}{\operatorname{def}}
\newcommand{\odd}{\operatorname{odd}}
\newcommand{\iso}{\operatorname{iso}}
\newcommand{\cN}{\mathcal{N}}
\newcommand{\cO}{\mathcal{O}}
\newcommand{\cE}{\mathcal{E}}
\newcommand{\cK}{\mathcal{K}}
\newcommand{\cP}{\mathcal{P}}
\newcommand{\cQ}{\mathcal{Q}}

\newcommand{\footremember}[2]{%
  \footnote{#2}%
  \newcounter{#1}%
  \setcounter{#1}{\value{footnote}}%
}

\newcommand{\ag}[1]{\textcolor{purple}{(Anna: #1)}}
\newcommand{\sd}[1]{\textcolor{red}{(Sahar: #1)}}
\title{Every vertex of a minimal Tutte obstruction has neighbours in at least three components of the graph minus that set}
\author{Sahar Diskin, Anna Geisler}
\date{arxiv:2404.14020v4; CC BY 4.0}
\begin{document}
\maketitle
\noindent\textit{Editorial note.} All delivered mathematical text is the author's own, copied from the pinned author source sources/combinatorics-fl04-vspace-1-5cm-perfect-matching-in-p/pinned-original-source.tex; no statement or proof step was rewritten, reordered or removed.\par\medskip
\noindent\textit{Editorial note.} Preamble provenance: the source pins the publisher class scrartcl, which the delivery wrapper does not accept (supported wrapper classes are article and amsart), so this card uses sources/combinatorics-fl04-vspace-1-5cm-perfect-matching-in-p/retained-preamble.tex, which supplies the amsart baseline and then carries the author's own newtheorem and macro declarations (pinned source lines 30-57) verbatim. No macro was invented and no mathematical symbol was changed.\par\medskip
\noindent\textit{Editorial note.} Bibliography provenance: the source itself carries no thebibliography environment and ends with bibliographystyle abbrv plus bibliography perc (author source lines 1691-1692), so the bibliography is resolved with BibTeX; the database used is the identical copy staged with the same arXiv version under sources/product-graph-matchings-2404.14020v4/perc.bib, whose licence evidence is also status supported for source\_id arxiv:2404.14020v4. No entry was edited and the bibliography style is the author's own abbrv.\par\medskip
\noindent\textit{Editorial note.} The displayed Tutte-violation relation that the counted statement and proof invoke, source label eq:tutte-violation, is delivered verbatim as a context span so that the reference in the counted text points at a formula present in the assembled material; it is the author's own display.\par\medskip
\noindent\textit{Editorial note.} The proof cites nothing and needs no navigation stub.\par\medskip
\noindent\textit{Allowed standard prerequisites (editorial).}\begin{itemize}
\item The Tutte-violation relation for a set U and the notion of a minimal set satisfying it. Stated by the author immediately before the counted lemma (the display labelled eq:tutte-violation at source line 660, delivered verbatim as a context span of this card). Fixes the property that the minimal set U has more components in G\_p - U than its own size, which the proof uses to count components after deleting a vertex.
\end{itemize}\medskip
\begin{lemma} \label{lem:three_edges}
Let $U$ be a minimal set of vertices satisfying \eqref{eq:tutte-violation}. Then every vertex $v \in U$ has neighbours (in $G$) in at least three components of $G_p-U$.
\end{lemma}

\begin{align} \label{eq:tutte-violation}
    \odd(G_p - U) \geq |U|+2
\end{align}

\begin{proof}
Suppose towards contradiction that there is a minimal set $U$ of size $u$ satisfying \eqref{eq:tutte-violation}, with $\ell>u$ components in $G_p-U$ and containing a vertex $v$ whose neighbourhood in $G$ intersects at most two components of $G_p-U$.

Consider $U'=U \setminus \{v\}$. Clearly, $|U'|=u-1$. Observe that since $v$ is connected to at most two components in $G_p-U$, there are at least $\ell-1>u-1=|U'|$ components in $G_p-U'$. But then $U'$ is an obstruction, contradicting the minimality of $U$.
\end{proof}

\bibliographystyle{plain}
\bibliography{perc}
\end{document}
